\documentclass[11pt,a4paper]{article}

\usepackage[T1]{fontenc}
\usepackage[utf8]{inputenc}
\usepackage[slovak]{babel}
\usepackage{lmodern}
\usepackage[margin=25mm]{geometry}
\usepackage{graphicx}
\usepackage{xcolor}
\usepackage{hyperref}

\hypersetup{colorlinks=true, linkcolor=black, urlcolor=blue!60!black}
\newcommand{\doplnit}[1]{\par\noindent\textit{Doplniť: #1}\par}
\newcommand{\obrazok}[3]{%
  \IfFileExists{#1}{%
    \begin{figure}[htbp]
      \centering
      \includegraphics[width=\linewidth,height=.72\textheight,keepaspectratio]{#1}
      \caption{#2}\label{#3}
    \end{figure}%
  }{\doplnit{Pridať obrázok \texttt{#1}: #2.}}%
}

\begin{document}

\hypersetup{pageanchor=false}
\begin{titlepage}
  \centering
  \includegraphics[width=0.28\textwidth]{assets/fiit-logo.png}\par
  \vspace{0.5cm}
  {\large Slovenská technická univerzita v Bratislave\par}
  {\large Fakulta informatiky a informačných technológií\par}

  \vfill
  {\Huge\bfseries Projektová dokumentácia\par}
  \vspace{0.7cm}
  {\Large VPWA --- Vývoj progresívnych webových aplikácií\par}
  \vspace{0.4cm}
  {\large Projekt: Convo --- aplikácia na textovú komunikáciu\par}
  \vfill

  \begin{tabular}{r l}
    \textbf{Študenti:} & Dávid Boldog, Tomáš Bubeník \\
    \textbf{Študijná skupina:} & Štvrtok 10:00 \\
    \textbf{Cvičiaci:} & Ing. Matúš Krajčovič \\
    \textbf{Akademický rok:} & 2026/2027 \\
    \textbf{Dátum:} & \today
  \end{tabular}
  \vspace{1cm}
\end{titlepage}

\hypersetup{pageanchor=true}
\pagenumbering{roman}
\tableofcontents
\clearpage
\pagenumbering{arabic}

\section{Zadanie}
\subsection{Cieľ projektu}
Convo je progresívna webová aplikácia na textovú komunikáciu v kanáloch.
Vychádza zo zadania semestrálneho projektu VPWA~\cite{vpwa-zadanie}.
\doplnit{Stručne opísať výslednú aplikáciu a jej cieľovú skupinu.}

\subsection{Prípady použitia}
Zadanie vyžaduje nasledujúcich jedenásť prípadov použitia:
\begin{enumerate}
  \item Registrácia, prihlásenie a odhlásenie používateľa.
  \item Zobrazenie, vytvorenie, opustenie a zrušenie kanálov.
  \item Odosielanie správ a príkazov cez pevne umiestnený vstup.
  \item Správa verejných a súkromných kanálov pomocou príkazov vrátane
        pozývania, odoberania členov a hlasovania o vylúčení.
  \item Zrušenie vlastného členstva; odchod správcu ukončí kanál.
  \item Adresovanie správy používateľovi pomocou zmienky.
  \item Prezeranie úplnej histórie správ s postupným načítavaním.
  \item Upozornenia na nové správy a voľba upozornení len na zmienky.
  \item Zmena stavu online, nerušiť a offline.
  \item Zobrazenie členov kanála.
  \item Zobrazenie používateľov, ktorí píšu, a náhľad ich rozpracovaného textu.
\end{enumerate}
\doplnit{Pri finálnej kontrole zosúladiť opis každého prípadu s implementáciou.}

\subsection{Použité technológie}
\doplnit{Uviesť výsledné verzie Quasaru, AdonisJS, relačnej databázy a
ďalších základných nástrojov.}

\section{Architektúra aplikácie}
\subsection{Diagram architektúry}
\obrazok{assets/architektura.png}{Architektúra aplikácie}{fig:architektura}
\doplnit{Vysvetliť tok požiadavky od prehliadača cez server až po databázu.
V diagrame označiť HTTP a komunikáciu v reálnom čase. Ak jeden diagram
nestačí, pridať ďalší.}

\subsection{Klientská časť}
\doplnit{Opísať stránky, hlavné komponenty, smerovanie, správu stavu,
responzivitu a napojenie Quasar aplikácie na server.}

\subsection{Serverová časť}
\doplnit{Opísať rozdelenie AdonisJS aplikácie na routy, kontroléry, služby
a modely; vysvetliť autentifikáciu a kontrolu oprávnení.}

\subsection{Databáza a perzistencia}
\doplnit{Vysvetliť, ktoré údaje sa ukladajú, ako server používa modely a
migrácie a ktoré krátkodobé stavy zostávajú iba v pamäti.}

\subsection{Komunikácia medzi klientom a serverom}
\doplnit{Oddeliť použitie HTTP API a spojenia v reálnom čase. Uviesť hlavné
udalosti pre správy, prítomnosť, písanie a rozpracovaný text.}

\section{Fyzický dátový model}
\subsection{Diagram modelu}
\obrazok{assets/fyzicky-model.png}{Fyzický dátový model}{fig:fyzicky-model}
\doplnit{Diagram vytvoriť podľa skutočných migrácií vrátane kľúčov, väzieb
a kardinalít.}

\subsection{Tabuľky, väzby a obmedzenia}
\doplnit{Vysvetliť význam tabuliek, dôležité stĺpce, unikátne obmedzenia,
referenčnú integritu a pravidlá mazania údajov.}

\subsection{Zmeny oproti logickému návrhu}
\doplnit{Porovnať fyzický model s logickým modelom z prvej fázy. Každú zmenu
vecne zdôvodniť; ak zmeny nie sú, výslovne to uviesť.}

\section{Realizácia funkcií}
\subsection{Používatelia a prihlásenie}
\doplnit{Opísať registráciu, prihlasovanie, odhlasovanie, identitu a
užívateľské nastavenia.}

\subsection{Kanály, členstvo a oprávnenia}
\doplnit{Opísať typy kanálov, členstvo, pozvánky a príkazy
\texttt{/join}, \texttt{/invite}, \texttt{/revoke}, \texttt{/kick},
\texttt{/cancel}, \texttt{/quit} a \texttt{/list}. Uviesť oprávnenia správcu
a zánik neaktívneho kanála.}

\subsection{Správy, zmienky a história}
\doplnit{Vysvetliť odoslanie správy, zmienky, ukladanie a postupné načítanie
histórie vrátane správania pri veľkom počte správ.}

\subsection{Upozornenia a stav používateľa}
\doplnit{Opísať podmienky vystavenia upozornenia, nastavenie iba pre zmienky
a účinok stavov online, nerušiť a offline.}

\subsection{Indikácia písania}
\doplnit{Vysvetliť zobrazenie píšucich členov a živý náhľad rozpracovaného
textu, vrátane ukončenia alebo vypršania tohto stavu.}

\section{PWA}
\subsection{Inštalácia a manifest}
\doplnit{Popísať manifest, ikony, režim zobrazenia a spôsob inštalácie
aplikácie na podporovaných zariadeniach.}

\subsection{Service worker a práca bez spojenia}
\doplnit{Vysvetliť stratégiu ukladania súborov, dostupnosť aplikácie bez
siete, správanie pri strate spojenia a opätovnú synchronizáciu.}

\subsection{Aktualizácie aplikácie}
\doplnit{Opísať, ako klient zistí a použije novú verziu aplikácie.}

\section{Návrhové rozhodnutia}
\subsection{Základné technológie a externé knižnice}
\doplnit{Vymenovať použité knižnice mimo základného rámca a pri každej
vysvetliť účel a dôvod výberu. Zdôvodniť aj prípadné úpravy cudzej knižnice.}

\subsection{Dátové a bezpečnostné pravidlá}
\doplnit{Zdôvodniť model členstva, pozvánok, hlasov, banov, správcovských
oprávnení a prístup k správam.}

\subsection{Výkon a správanie v reálnom čase}
\doplnit{Zdôvodniť stránkovanie histórie, prenos udalostí, obnovu spojenia
a spôsob aktualizácie klientského stavu.}

\subsection{Obmedzenia riešenia}
\doplnit{Uviesť iba skutočné obmedzenia finálnej verzie a ich dopad na
používateľa. Ak žiadne významné nie sú, túto podkapitolu odstrániť.}

\section{Kľúčové obrazovky}
Vložiť najmenej päť snímok finálnej aplikácie. Pri každej uviesť stručný
popis toho, čo používateľ na obrazovke robí. Súbory uložiť do priečinka
\texttt{assets/}.

\subsection{Prihlásenie a registrácia}
\obrazok{assets/prihlasenie.png}{Prihlásenie a registrácia}{fig:prihlasenie}
\doplnit{Popísať scenár na snímke.}

\subsection{Kanály a pozvánky}
\obrazok{assets/kanaly.png}{Kanály a pozvánky}{fig:kanaly}
\doplnit{Popísať scenár na snímke.}

\subsection{Konverzácia a história správ}
\obrazok{assets/konverzacia.png}{Konverzácia a história správ}{fig:konverzacia}
\doplnit{Popísať scenár na snímke.}

\subsection{Správa členov kanála}
\obrazok{assets/clenovia.png}{Správa členov kanála}{fig:clenovia}
\doplnit{Popísať scenár na snímke.}

\subsection{Nastavenia, stav a upozornenia}
\obrazok{assets/nastavenia.png}{Nastavenia, stav a upozornenia}{fig:nastavenia}
\doplnit{Popísať scenár na snímke.}

\section{Spustenie a overenie}
\subsection{Požiadavky a konfigurácia}
\doplnit{Uviesť potrebné verzie nástrojov, príklad premenných prostredia
bez tajných hodnôt a požiadavky na databázu.}

\subsection{Inštalácia a spustenie}
\doplnit{Zapísať presné príkazy na inštaláciu balíkov, migráciu databázy
a spustenie klienta a servera.}

\subsection{Overenie prípadov použitia}
\doplnit{Stručne popísať, ako ste overili všetkých jedenásť prípadov
použitia, vrátane dvoch súčasne prihlásených používateľov.}

\subsection{Overenie PWA}
\doplnit{Uviesť výsledok skúšky inštalácie, režimu bez spojenia, návratu
online a aktualizácie aplikácie.}

\begin{thebibliography}{9}
  \bibitem{vpwa-zadanie}
  Zadanie semestrálneho projektu VPWA.\newline
  \url{https://github.com/kurice/vpwa27/blob/main/semestralny-projekt/readme.md}

\end{thebibliography}

\end{document}
